\documentclass[10pt,a4paper]{article}

% Packages
\usepackage[utf8]{inputenc}
\usepackage[T1]{fontenc}
\usepackage{geometry}
\usepackage{titlesec}
\usepackage{enumitem}
\usepackage{hyperref}
\usepackage{xcolor}
\usepackage{fontawesome5}
\usepackage{helvet} % Use Helvetica-like font

% Formatting
\renewcommand{\familydefault}{\sfdefault} % Set default font to Sans-Serif
\pagestyle{empty} % No page numbers
\setlength{\parindent}{0pt} % No indentation

% Section Formatting
\titleformat{\section}{
  \vspace{-4pt}\scshape\raggedright\large\bfseries
}{}{0em}{}[\color{black}\titlerule \vspace{-4pt}]

% Custom commands
\newcommand{\resumeItem}[1]{
  \item\small{
    {#1 \vspace{-2pt}}
  }
}

\newcommand{\resumeSubheading}[4]{
  \vspace{-1pt}\item
    \begin{tabular*}{0.97\textwidth}[t]{l@{\extracolsep{\fill}}r}
      \textbf{#1} & #2 \\
      \textbf{#3} & #4 \\ % Changed to Bold for Role/Degree based on image impression, or keep as is? Let's format as Title (Bold) + details
    \end{tabular*}\vspace{-5pt}
}

% Redefining to match the image: 
% The image has:
% EMPLYOYER ................. LOCATION
% ROLE ...................... DATE
% Both Employer and Role look prominent. 
% Actually, looking at the image: "Bachelor of Informatics" is Bold. "Marketing Campaign Analyst" is NOT Bold.
% So:
% School: Bold
% Degree: Bold
% Employer: Bold
% Job Title: Normal?
% Let's try to make it flexible or just follow the standard "Jake's Resume" which is usually:
% Bold School, Italic Degree. 
% The user said "layout nya seperti gambar ini". The image uses Sans Serif. 
% Let's stick to ResumeSubheading being: #1 Bold, #2 Normal, #3 Italic/Normal, #4 Italic/Normal.
% But I will switch to Sans Serif which is the biggest change.

\renewcommand{\resumeSubheading}[4]{
  \vspace{2pt}\item
    \begin{tabular*}{0.97\textwidth}[t]{l@{\extracolsep{\fill}}r}
      \textbf{#1} & #2 \\
      #3 & #4 \\
    \end{tabular*}\vspace{-4pt}
}

\newcommand{\resumeProjectHeading}[2]{
    \vspace{2pt}\item
    \begin{tabular*}{0.97\textwidth}{l@{\extracolsep{\fill}}r}
      \textbf{#1} & #2 \\
    \end{tabular*}\vspace{-4pt}
}

\newcommand{\resumeSubItem}[1]{\resumeItem{#1}\vspace{-4pt}}

\renewcommand\labelitemii{$\vcenter{\hbox{\tiny$\bullet$}}$}

\newcommand{\resumeSubHeadingListStart}{\begin{itemize}[leftmargin=0.0in, label={}]} % Changed margin to 0 to align with section headers
\newcommand{\resumeSubHeadingListEnd}{\end{itemize}}
\newcommand{\resumeItemListStart}{\begin{itemize}[leftmargin=0.2in]} % Adjusted bullet indentation
\newcommand{\resumeItemListEnd}{\end{itemize}\vspace{-4pt}}

\begin{document}
\begin{center}
    \textbf{\Huge \scshape Rizky Yanuar Kristianto} \\ \vspace{5pt}
    Surabaya, Jawa Timur, Indonesia \\
    Phone: +6289676257683 \\
    E-mail: \href{mailto:rizkyyanuarkristianto@gmail.com}{rizkyyanuarkristianto@gmail.com} \\
    \href{https://linkedin.com/in/rizkyyanuark}{www.linkedin.com/in/rizkyyanuark} \\
    \href{https://github.com/rizkyyanuark}{https://github.com/rizkyyanuark}
\end{center}

\vspace{-10pt}
\begin{center}
    \small \parbox{0.95\textwidth}{
        Mahasiswa Sains Data semester akhir dengan pengalaman magang sebagai \textbf{AI Engineer} dan \textbf{Data Engineer}. Terampil membangun solusi AI/ML end-to-end, mulai dari ETL pipeline hingga chatbot berbasis NLP. Siap berkontribusi dalam transformasi digital dan inovasi berbasis data.
    }
\end{center}

%-----------EDUCATION-----------
\section{Pendidikan}
\resumeSubHeadingListStart
\resumeSubheading
{Universitas Negeri Surabaya}{Surabaya, Jawa Timur}
{\textbf{Sarjana Sains Data (S1)}}{Agu 2022 -- Sekarang}
\resumeItemListStart
\resumeItem{\textbf{IPK}: 3.53/4.00 (103 SKS)}
\resumeItem{Mata Kuliah Relevan: Sains Data, Struktur Data dan Algoritma, Statistika, Kecerdasan Buatan (AI), Machine Learning, Computer Vision, Generative AI.}
\resumeItemListEnd
\resumeSubHeadingListEnd

%-----------EXPERIENCE-----------
\section{Pengalaman Kerja}
\resumeSubHeadingListStart

\resumeSubheading
{Dinas Komunikasi dan Informatika Provinsi Jawa Timur}{Surabaya, Jawa Timur}
{\textbf{AI Engineer Intern}}{Sep 2024 -- Jan 2025}
\resumeItemListStart
\resumeItem{Terlibat dalam program \textbf{Percepatan 100 Hari Kerja Gubernur Jawa Timur} bersama tim pengembangan layanan digital.}
\resumeItem{Membantu pengembangan fitur chatbot berbasis AI untuk aplikasi \textbf{Majadigi}, melayani \textbf{1.000+} query harian.}
\resumeItem{Mengimplementasikan komponen \textbf{NLP} untuk memproses pertanyaan pengguna secara otomatis.}
\resumeItem{Menyusun dokumentasi teknis dan membantu pelatihan \textbf{5} anggota tim operasional.}
\resumeItemListEnd

\resumeSubheading
{PT PLN Icon Plus}{Surabaya, Jawa Timur}
{\textbf{Data Engineer Intern}}{Feb 2024 -- Jun 2024}
\resumeItemListStart
\resumeItem{Membangun \textbf{ETL pipeline} menggunakan Apache Airflow dan PySpark untuk mengotomatisasi pengolahan \textbf{50.000+} data operasional harian.}
\resumeItem{Mengembangkan prototipe chatbot berbasis \textbf{RAG} untuk \textbf{20+} staf, mempercepat akses informasi internal.}
\resumeItem{Melakukan integrasi data dari \textbf{3} sumber ke dalam \textbf{Data Warehouse} PostgreSQL untuk kebutuhan pelaporan.}
\resumeItem{Berkolaborasi dengan tim dalam iterasi perbaikan model AI berdasarkan feedback pengguna.}
\resumeItemListEnd

\resumeSubheading
{Universitas Negeri Surabaya}{Surabaya, Jawa Timur}
{\textbf{Asisten Peneliti (Ketua Tim Mahasiswa)}}{Jun 2023 -- Des 2023}
\resumeItemListStart
\resumeItem{Memimpin tim \textbf{4} mahasiswa dalam proyek riset deteksi depresi berbasis analisis suara.}
\resumeItem{Membangun model \textbf{CNN-LSTM} dengan TensorFlow pada \textbf{800+} sampel audio, mencapai akurasi \textbf{85\%}.}
\resumeItem{Berkontribusi dalam penulisan laporan riset dan pendaftaran \textbf{HKI}.}
\resumeItemListEnd

\resumeSubHeadingListEnd

%-----------ORGANIZATIONAL EXPERIENCE-----------
\section{Pengalaman Organisasi}
\resumeSubHeadingListStart
\resumeSubheading
{Himpunan Mahasiswa Sains Data}{Surabaya, Jawa Timur}
{\textbf{Staf Divisi Advokasi dan Kesejahteraan Mahasiswa}}{Jan 2024 -- Jan 2025}
\resumeItemListStart
\resumeItem{Membantu penyelenggaraan program \textbf{InSCAW} (Career Workshop) untuk pengembangan karir mahasiswa.}
\resumeItem{Menjadi penghubung antara mahasiswa dan fakultas dalam koordinasi kegiatan akademik.}
\resumeItemListEnd
\resumeSubHeadingListEnd

%-----------PROJECTS-----------
\section{Proyek}
\resumeSubHeadingListStart
\resumeProjectHeading
{\textbf{Academic Search Assistant with Knowledge Graph RAG (Tugas Akhir)}}{Jan 2025 -- Saat Ini}
\resumeItemListStart
\resumeItem{Mengembangkan sistem pencarian literatur akademik menggunakan pendekatan \textbf{Hybrid Vector-Graph RAG}.}
\resumeItem{Membangun \textbf{Knowledge Graph} untuk meningkatkan konteks dan keterlacakan referensi dalam respons LLM.}
\resumeItemListEnd
\resumeProjectHeading
{\textbf{NV-Bite: Android App with Cloud \& ML Integration (Bangkit Academy)}}{Agu 2024 -- Des 2024}
\resumeItemListStart
\resumeItem{Mengikuti program \textbf{Bangkit Academy 2024} dan meraih penghargaan \textbf{Best Presenting Group} dalam presentasi capstone project.}
\resumeItem{Berkontribusi dalam pengembangan model klasifikasi gambar menggunakan TensorFlow dan deployment di \textbf{Google Cloud Platform}.}
\resumeItemListEnd
\resumeProjectHeading
{\textbf{KIPK Scholarship Sentiment Analysis}}{Jun 2024}
\resumeItemListStart
\resumeItem{Membangun pipeline analisis sentimen pada data Twitter menggunakan model \textbf{BiLSTM} dengan Python.}
\resumeItem{Membuat visualisasi hasil analisis dalam dashboard \textbf{Looker Studio} untuk presentasi kepada stakeholder.}
\resumeItemListEnd
\resumeProjectHeading
{\textbf{Smart Sudoku Solver dengan Computer Vision}}{Jan 2024 -- Agu 2024}
\resumeItemListStart
\resumeItem{Mengembangkan aplikasi penyelesai Sudoku otomatis menggunakan \textbf{OpenCV} untuk deteksi grid dan \textbf{CNN} untuk pengenalan digit.}
\resumeItem{Mengimplementasikan algoritma backtracking untuk menyelesaikan puzzle secara programatis.}
\resumeItemListEnd
\resumeSubHeadingListEnd

%-----------SKILLS-----------
\section{Keahlian}
\textbf{Keahlian Teknis:}
\begin{itemize}[leftmargin=0.2in, itemsep=-2pt, topsep=0pt]
    \small\item \textbf{Bahasa Pemrograman:} Python, SQL, Bash, R.
    \item \textbf{AI \& Machine Learning:} TensorFlow, PyTorch, Keras, Scikit-learn, Langchain (LLM), OpenCV (Computer Vision), NLP.
    \item \textbf{Data Engineering \& Cloud:} Apache Airflow, PySpark, PostgreSQL, MySQL, Google Cloud Platform (GCP), BigQuery, Docker.
    \item \textbf{Visualisasi Data:} Matplotlib, Looker Studio, Tableau, Excel.
    \item \textbf{Keahlian Non-Teknis:} Manajemen Proyek, Kepemimpinan Tim, Komunikasi Strategis, Pemecahan Masalah, Adaptabilitas.
\end{itemize}
\vspace{2pt} % Positive spacing to separate categories

\textbf{Sertifikasi \& Pelatihan:}
\begin{itemize}[leftmargin=0.2in, itemsep=-2pt, topsep=0pt]
    \small\item \textbf{JuaraGCP Session 11} -- Google Cloud
    \item \textbf{DeepLearning.AI TensorFlow Developer} -- DeepLearning.AI
    \item \textbf{TensorFlow: Data and Deployment} -- DeepLearning.AI
    \item \textbf{Generative AI for Everyone} -- DeepLearning.AI
    \item \textbf{Learn Data Analysis with Python} -- Dicoding
\end{itemize}
\vspace{2pt} % Positive spacing to separate categories

\textbf{Bahasa:}
\begin{itemize}[leftmargin=0.2in, itemsep=-2pt, topsep=0pt]
    \small\item \textbf{Fasih:} Bahasa Indonesia, Bahasa Inggris (\textbf{IELTS: 7.5})
    \item \textbf{Konversasional:} Bahasa Jawa
\end{itemize}

\end{document}
